\chapter{Inert-Prime Derivative Divisibility and Cyclic Orbital Linear Programming}
\label{ch:08}

In \cref{ch:07} we proved that at every level-free inert prime $p \ge 5$, the corrected Weil theta kernels along our definite unitary tower are $p$-adically integral (\cref{thm:07:int-inert}) and that the height-$2$ Lubin--Tate biextension satisfies $\beta(X, Y) \sim G_{c, \sigma(c)}(X, Y)$ (\cref{thm:07:biextension-rigidity}). We now exploit this height-$2$ structure to extract \emph{explicit high-order $p$-adic divisibility} from the formal-group derivatives of the theta kernels along cyclic Frobenius orbits of arbitrary even or odd length $d$.

Two fundamental structural engines power this chapter:
\begin{enumerate}[label=(\roman*)]
\item \emph{Height-$2$ pole-to-LP duality.} Expanding $\ell_p(X) = \sum_{k \ge 0} (-p)^{-k} X^{p^{2k}}$ and $G_{c,\sigma(c)}(X, Y)$ shows that the derivative of order $D(s)$ at Frobenius orbit position $s \in \ZZ/d\ZZ$ loses valuation only at odd powers $p^{2k+1}$ paired across distance $i \in \mathrm{oddLevels}(d)$, with exact asymptotic slope $\gamma_p = \frac{p}{p^2 - 1}$. Maximizing the achievable valuation gain over a cyclic orbit of length $d$ is therefore a finite-dimensional linear program!
\item \emph{Zero-gap Euclidean-projection LP duality and sharpened constant $8$.} We replace Hahn--Banach and Tychonoff compactness in the proof of strong LP duality (\cref{thm:08:lp-duality}) by a fourteen-line, zero-trap Bolzano--Weierstrass proper-cone Euclidean projection proof, and we collapse the odd-level perturbation analysis to an eight-line one-shot dual test (\cref{thm:08:inert-deriv}) that simultaneously sharpens the explicit loss constant from $11$ down to $8$!
\end{enumerate}

\section{Level-Free Inert Places, Frobenius Orbits, and Engine~I}
\label{sec:08:engine-one}

Let $\ell = p \ge 5$ be a level-free rational prime inert in $K = \QQ(\sqrt{-d})$, meaning $\ell \nmid 6 Q_S q d g_0 g_1 T$. Because $\ell$ is unramified in $E/\QQ$ and $\Frob_\ell \in G \setminus \Phi$, the Frobenius action $\tau \mapsto \Frob_\ell \cdot \tau$ partitions the CM embeddings $\Phi$ (and their conjugates $c\Phi$) into cyclic orbits of even length $2d$ (or $d$ on $\Phi$ under $\Frob_\ell^2$ when viewed on pairs $(\tau, c\tau)$). Write $\ell \ge 5$ throughout this chapter for the inert prime (matching the cyclic orbital LP notation $\mathcal{D}_{\ell, d}$), and define the \emph{height-$2$ canonical slope}
\begin{equation}
\label{eq:08:gamma-p}
  \gamma_\ell \;:=\; \frac{\ell}{\ell^2 - 1} \;\in\; \left(\frac{1}{\ell},\; \frac{3}{2\ell}\right] \qquad (\ell \ge 3).
\end{equation}
Notice why $\gamma_\ell = \frac{\ell}{\ell^2 - 1}$ governs height-$2$ formal groups: in the Lubin--Tate logarithm $\ell_\ell(X) = \sum_{k \ge 0} (-\ell)^{-k} X^{\ell^{2k}}$, differentiating $D$ times produces a denominator of $\ell$-adic valuation $\approx \lfloor D / \ell^2 \rfloor + \lfloor D / \ell^4 \rfloor + \cdots \approx \frac{D}{\ell^2 - 1}$ compared to ordinary Legendre factorial valuation $\frac{D}{\ell - 1}$; the difference $\frac{1}{\ell - 1} - \frac{1}{\ell^2 - 1} = \frac{\ell}{\ell^2 - 1} = \gamma_\ell$ is the exact valuation excess harvested by height-$2$ supersingular reduction!

\begin{definition}[Level-free inert place data and true-column stabilization]
\label{def:08:level-free}
A rational prime $\ell \ge 5$ inert in $K$ is \emph{level-free} if $\ell$ is unramified in $E/\QQ$ and coprime to all auxiliary level and supercuspidal primes $g_0, g_1, T$. At any level-free inert prime $\ell$:
\begin{enumerate}[label=\textup{(\roman*)}]
\item The local host realization exists without ideal coprimality restrictions (by moving $\mathfrak{a}_0$ inside its genus to a $\ell$-integral representative via weak approximation), giving local $\OO_{K_\ell}$-orthonormal bases at every tower layer $0 \le j \le m$.
\item Evaluating a column chain (\cref{def:07:tower-ranks}) on the formal neighborhood of CM torsion points via the corrected theta kernel $\theta_{MM}$ of \cref{thm:07:int-inert} stabilizes onto \emph{true columns}: every entry in degree vector $D = (D(s))_{s \in \ZZ/d\ZZ}$ lies in $\OO_{\CC_\ell}$ up to the universal Hurwitz factor determined by the Lubin--Tate biextension $G_{c,\sigma(c)}$ of \cref{thm:07:biextension-rigidity}.
\end{enumerate}
\end{definition}
\leanref{src/Formalism/SharpRestriction/Divisibility/OrbitPlaces.lean:64}{Formalism.MainNotions.OrbitPlaces.LevelFree}
\leanref{src/Development/AbelianVarietyTheory/Proofs/SatakeParameters/JacobiOfClass.lean:63}{GlobalTheory.KerStd.HostInert.hostAt\_inert}
\leanref{src/Development/AbelianVarietyTheory/Proofs/SatakeParameters/AtOfTrueColumns.lean:30}{GlobalTheory.PdvClose.at\_of\_true\_columns}

We first record \emph{Engine~I}: the direct Frobenius-lift recursive derivative valuation bound obtained by differentiating $\exp(G_{c,\sigma(c)}(X, Y))$ along the $\ell$-power Frobenius filtration without invoking any auxiliary operator algebra.

\begin{proposition}[Engine~I: Generic Frobenius-lift kernel derivative valuation bound]
\label{prop:08:engine-one}
Let $\ell \ge 5$ be a level-free inert prime, let $\mathcal{O}$ be an orbit of length $d$ under $\Frob_\ell$ with nonnegative integer derivative-order profile $D : \ZZ/d\ZZ \to \NN$, and let $\mathrm{oddLevels}(d) := \{i \in \{1, 3, \dots, d-1\} : i \text{ odd}\}$ (when $d$ is even; for odd $d$, unfold to period $2d$). For any nonnegative flow $x : \ZZ/d\ZZ \times \mathrm{oddLevels}(d) \to \RR_{\ge 0}$ satisfying the \emph{orbital degree constraint}
\begin{equation}
\label{eq:08:flow-constraint}
  (A x)(t) \;:=\; \sum_{i \in \mathrm{oddLevels}(d)} \ell^i\, x(t, i) \;+\; \sum_{i \in \mathrm{oddLevels}(d)} x(t - i, i) \;\le\; D(t)
  \qquad \forall\, t \in \ZZ/d\ZZ,
\end{equation}
the normalized $p$-adic valuation of every formal-group derivative of degree profile $D$ of the corrected theta kernel $\theta_{MM}$ satisfies the recursive orbit gain bound
\begin{equation}
\label{eq:08:engine-one-bound}
  v_\ell\bigl(\partial^D \theta_{MM}\bigr)
  \;\ge\; \gamma_\ell \sum_{s \in \ZZ/d\ZZ} D(s) \;-\; \mathbf{1}^T x \;-\; \mathrm{Err}_{\mathrm{I}}(\ell, d, D),
\end{equation}
where $\mathbf{1}^T x := \sum_{s, i} x(s, i)$ is the total flow value and $\mathrm{Err}_{\mathrm{I}}$ is bounded linearly in $(5, 4)$ by the generic boundary rounding terms (specifically at most $5 d + 4 \sum_s \lfloor D(s) / \ell^{d} \rfloor$). In particular, taking the supremum over all feasible flows $x$ yields the \emph{primal orbital LP value}
\[
  v_{\mathrm{primal}}(\ell, d, D) \;:=\; \sup\bigl\{\mathbf{1}^T x : x \in \RR_{\ge 0}^{E},\; A x \le D\bigr\},
  \qquad E := \ZZ/d\ZZ \times \mathrm{oddLevels}(d),
\]
and defines the \emph{orbital gain functional}
\begin{equation}
\label{eq:08:fnum-def}
  F_{\mathrm{num}}(\ell, d, D) \;:=\; \gamma_\ell \sum_{s \in \ZZ/d\ZZ} D(s) \;-\; v_{\mathrm{primal}}(\ell, d, D).
\end{equation}
\end{proposition}
\leanref{src/Formalism/ThetaLogCoordinates/UniformThetaSubstitution/FrobeniusLiftNorms/EGaussK.lean:255}{Formalism.InvariantFactorPaths.G7.kernel\_bounds\_generic}
\leanref{src/Development/AbelianVarietyTheory/Proofs/SatakeParameters/KvalGeOrbitGainWithoutOperator.lean:106}{GlobalTheory.KerGaussA.kvalGe\_gainA}

\begin{proof}
By \cref{thm:07:int-inert,thm:07:biextension-rigidity}, the corrected theta kernel expands on a $p$-local orthonormal basis as an $\OO_{\CC_\ell}$-linear combination of Hurwitz-integral power series multiplied by $\exp(G_{c, \sigma(c)}(X_s, X_{s'}))$ along Frobenius-conjugate coordinate pairs $(X_s, X_{s + i})$ with $i \in \mathrm{oddLevels}(d)$. Differentiating $\exp(\sum c_{s,i} \ell^{-k} X_s^{\ell^i} X_{s+i})$ to multidegree $D = (D(s))_{s \in \ZZ/d\ZZ}$ expands as a sum over all monomial factorizations allocating $x(s, i)$ copies of the singular term $\ell^{-1} X_s^{\ell^i} X_{s+i}$ (consuming degree $\ell^i x(s, i)$ at position $s$ and degree $x(s, i)$ at position $s + i$, while incurring valuation loss $1$ per copy!).

Summing the degree consumption at each position $t \in \ZZ/d\ZZ$ gives precisely the inequality $(A x)(t) = \sum_i \ell^i x(t, i) + \sum_i x(t - i, i) \le D(t)$ in \eqref{eq:08:flow-constraint}. Meanwhile, the remaining unallocated derivatives come from the Hurwitz-integral factor and the Lubin--Tate exponential $\exp(\ell_\ell(X_s))$, whose normalized Taylor coefficient of $X_s^m$ carries $\ell$-adic valuation at least $\gamma_\ell m - O(1)$ (by Legendre's formula $v_\ell(m!) = \frac{m - s_\ell(m)}{\ell - 1}$ minus the Lubin--Tate pole valuation $\frac{m - s_{\ell^2}(m)}{\ell^2 - 1}$). Taking the minimum over all integral allocations $x$ (which is bounded below by replacing the integer maximum $\max \mathbf{1}^T x$ by the continuous LP relaxation $v_{\mathrm{primal}}(\ell, d, D)$) yields \eqref{eq:08:engine-one-bound} with explicit rounding constants $(5, 4)$ from the finite number of odd-level carry terms.
\end{proof}

\section{Fourteen-Line Proper-Cone Euclidean-Projection Proof of Cyclic LP Strong Duality}
\label{sec:08:lp-duality}

To bound $F_{\mathrm{num}}(\ell, d, D) = \gamma_\ell \sum_s D(s) - v_{\mathrm{primal}}(\ell, d, D)$ from \emph{below}—which requires bounding the primal maximization LP $v_{\mathrm{primal}}(\ell, d, D)$ from \emph{above}—we pass to the dual minimization LP over the dual polytope
\begin{equation}
\label{eq:08:dual-polytope}
  \mathcal{D}_{\ell, d} \;:=\; \bigl\{\sigma : \ZZ/d\ZZ \to [0, 1] \;\big|\; \ell^i \sigma(s) + \sigma(s + i) \ge 1\ \ \forall\, (s, i) \in \ZZ/d\ZZ \times \mathrm{oddLevels}(d)\bigr\}.
\end{equation}
Notice that for any feasible primal flow $x \in \RR_{\ge 0}^E$ ($A x \le D$) and any feasible dual vector $\sigma \in \mathcal{D}_{\ell, d}$, \emph{weak duality} is a one-line algebraic identity:
\begin{equation}
\label{eq:08:weak-duality}
  \mathbf{1}^T x \;=\; \sum_{s, i} x(s, i) \;\le\; \sum_{s, i} \bigl(\ell^i \sigma(s) + \sigma(s + i)\bigr) x(s, i)
  \;=\; \sum_{t \in \ZZ/d\ZZ} \sigma(t)\,(A x)(t) \;\le\; \langle \sigma, D \rangle.
\end{equation}
Hence $v_{\mathrm{primal}}(\ell, d, D) \le \inf_{\sigma \in \mathcal{D}_{\ell, d}} \langle \sigma, D \rangle$ and $F_{\mathrm{num}}(\ell, d, D) \ge \sup_{\sigma \in \mathcal{D}_{\ell, d}} \sum_s D(s)(\gamma_\ell - \sigma(s))$.

To perturb between the true degree profile $D_\ell$ and the smooth base profile $W = (k - 1)w$, we also need \emph{strong duality}: the infimum is attained at an optimal $\sigma^* \in \mathcal{D}_{\ell, d}$ with $\langle \sigma^*, D \rangle = v_{\mathrm{primal}}(\ell, d, D)$. Rather than invoking infinite-dimensional Hahn--Banach separation or Farkas--Minkowski--Weyl polyhedral theory, we give a completely self-contained, fourteen-line proof using nothing beyond Bolzano--Weierstrass and Euclidean closest-point projection in $\RR^{d+1}$! Crucially, because every column of $A$ has entry $\ell^i \ge 1$ in row $s$, the cone $K_{\mathrm{set}}$ is \emph{proper} and proved closed in three lines with zero upper-boundary traps.

\newsimplification{Fourteen-Line Proper-Cone Euclidean-Projection Proof of Cyclic Orbital LP Strong Duality}

\begin{theorem}[Strong LP duality for arbitrary cyclic Frobenius orbits]
\label{thm:08:lp-duality}
Let $\ell \ge 2$ and $d \ge 2$ be integers with $\mathrm{oddLevels}(d) \neq \varnothing$, let $E := \ZZ/d\ZZ \times \mathrm{oddLevels}(d)$, and let $D : \ZZ/d\ZZ \to \RR_{\ge 0}$ be any nonnegative degree profile. Then both the primal maximum $v := \sup\{\mathbf{1}^T x : x \in \RR_{\ge 0}^E,\; A x \le D\}$ and the dual minimum $\inf_{\sigma \in \mathcal{D}_{\ell, d}} \langle \sigma, D \rangle$ are attained and strictly equal:
\begin{equation}
\label{eq:08:strong-duality}
  \max\bigl\{\mathbf{1}^T x : x \in \RR_{\ge 0}^E,\; A x \le D\bigr\}
  \;=\; \min_{\sigma \in \mathcal{D}_{\ell, d}} \langle \sigma, D \rangle.
\end{equation}
\end{theorem}
\leanref{src/Formalism/IncomingBounds/FlowDegreeBounds/MaxFlowValueAttained.lean:430}{Formalism.IncomingBounds.strongDuality}

\begin{proof}
1. \emph{Primal attainment.} For any feasible $x \in \RR_{\ge 0}^E$ ($A x \le D$) and any $(s, i) \in E$, we have $0 \le x(s, i) \le \ell^i x(s, i) \le (A x)(s) \le D(s)$. Thus the feasible set is nonempty ($0$ is feasible) and compact in $\RR^E$, so $v \ge 0$ is attained at some optimal flow $x^*$.

2. \emph{Cone definition.} Define the convex cone $K_{\mathrm{set}} \subseteq \RR^d \times \RR$ by
\[
  K_{\mathrm{set}} \;:=\; \bigl\{(q_1, q_2) \in \RR^d \times \RR \;\big|\; \exists\, x \in \RR_{\ge 0}^E,\;\; A x \le q_1 \text{ and } q_2 \le \mathbf{1}^T x\bigr\}.
\]
Clearly $0 \in K_{\mathrm{set}}$, and $K_{\mathrm{set}}$ is convex and closed under multiplication by scalars $\lambda \ge 0$.

3. \emph{Closedness of $K_{\mathrm{set}}$ in three lines.} Let $(q_1^{(n)}, q_2^{(n)}) \in K_{\mathrm{set}}$ converge to $(p_1, p_2) \in \RR^d \times \RR$ with witnesses $x^{(n)} \in \RR_{\ge 0}^E$. Set $R := \sup_n \|q_1^{(n)}\|_\infty < \infty$. Because $\ell^i \ge 1$ for every $(s, i) \in E$,
\[
  0 \;\le\; x^{(n)}(s, i) \;\le\; \ell^i x^{(n)}(s, i) \;\le\; (A x^{(n)})(s) \;\le\; q_1^{(n)}(s) \;\le\; R \qquad \forall\, n,\; \forall\, (s, i) \in E.
\]
Hence $\{x^{(n)}\}$ lies in the compact cube $[0, R]^E$. By Bolzano--Weierstrass, a subsequence $x^{(n_k)}$ converges to some $x_\infty \in [0, R]^E$; passing to the limit in $A x^{(n_k)} \le q_1^{(n_k)}$ and $q_2^{(n_k)} \le \mathbf{1}^T x^{(n_k)}$ gives $(p_1, p_2) \in K_{\mathrm{set}}$. Thus $K_{\mathrm{set}}$ is closed!

4. \emph{Euclidean closest-point projection.} Fix any $\varepsilon > 0$ and put $p_\varepsilon := (D, v + \varepsilon) \in \RR^d \times \RR$. By definition of $v$, $p_\varepsilon \notin K_{\mathrm{set}}$. Because $0 \in K_{\mathrm{set}}$ and $K_{\mathrm{set}}$ is closed, the intersection $K_{\mathrm{set}} \cap \overline{B}(p_\varepsilon, \|p_\varepsilon\|_2)$ is nonempty and compact in $\RR^{d+1}$, so $q \mapsto \|q - p_\varepsilon\|_2^2$ attains its global minimum on $K_{\mathrm{set}}$ at some point $q^* \in K_{\mathrm{set}}$.

5. \emph{Cone separation.} Set $n = (a, c) := q^* - p_\varepsilon \neq 0$ ($a \in \RR^d$, $c \in \RR$). For every $q \in K_{\mathrm{set}}$ and $t \in [0, 1]$, convexity gives $(1-t)q^* + t q \in K_{\mathrm{set}}$, so
\[
  0 \;\le\; \|(q^* - p_\varepsilon) + t(q - q^*)\|_2^2 - \|q^* - p_\varepsilon\|_2^2
  \;=\; 2t\,\langle n, q - q^* \rangle + t^2\|q - q^*\|_2^2.
\]
Dividing by $t \to 0^+$ yields $\langle n, q - q^* \rangle \ge 0$ for all $q \in K_{\mathrm{set}}$. Testing $q = 0 \in K_{\mathrm{set}}$ and $q = 2q^* \in K_{\mathrm{set}}$ forces $\langle n, q^* \rangle = 0$, whence
\[
  \langle n, q \rangle \;\ge\; 0 \quad (\forall\, q \in K_{\mathrm{set}}),
  \qquad\text{and}\qquad
  \langle n, p_\varepsilon \rangle \;=\; \langle n, q^* \rangle - \|n\|_2^2 \;=\; -\|n\|_2^2 \;<\; 0.
\]

6. \emph{Dual extraction.} Testing $q = (\mathbf{e}_t, 0) \in K_{\mathrm{set}}$ (witness $x = 0$) gives $a(t) \ge 0$ for all $t \in \ZZ/d\ZZ$. Testing $q = (A\mathbf{e}_{(s,i)}, 1) \in K_{\mathrm{set}}$ (witness $x = \mathbf{e}_{(s,i)}$) gives $\ell^i a(s) + a(s + i) + c \ge 0$ for all $(s, i) \in E$. Since $\langle a, D \rangle + c(v + \varepsilon) = \langle n, p_\varepsilon \rangle < 0$ while $a \ge 0$, $D \ge 0$, and $v + \varepsilon > 0$ (strictly positive even when $D \equiv 0$!), we must have $c < 0$.

7. \emph{Clipping and compactness.} Put $\tau_\varepsilon := a / (-c) \ge 0$, which satisfies $\ell^i \tau_\varepsilon(s) + \tau_\varepsilon(s + i) \ge 1$ for all $(s, i) \in E$ and $\langle \tau_\varepsilon, D \rangle < v + \varepsilon$. Clipping $\widetilde{\tau}_\varepsilon(s) := \min(\tau_\varepsilon(s), 1) \in [0, 1]$ preserves feasibility ($\ell^i \widetilde{\tau}_\varepsilon(s) + \widetilde{\tau}_\varepsilon(s + i) \ge 1$: if $\tau_\varepsilon(s) \ge 1$ then $\ell^i \cdot 1 \ge 1$; if $\tau_\varepsilon(s+i) \ge 1$ then $1 \ge 1$; otherwise neither coordinate was clipped) and obeys $\langle \widetilde{\tau}_\varepsilon, D \rangle \le \langle \tau_\varepsilon, D \rangle < v + \varepsilon$. Thus $\widetilde{\tau}_\varepsilon \in \mathcal{D}_{\ell, d} \subseteq [0, 1]^d$. Taking $\varepsilon = 1/m$ ($m \ge 1$), Bolzano--Weierstrass extracts a convergent subsequence $\widetilde{\tau}_{1/m_k} \to \sigma^* \in \mathcal{D}_{\ell, d}$ with $\langle \sigma^*, D \rangle \le v$. Combining with weak duality \eqref{eq:08:weak-duality} yields $\langle \sigma^*, D \rangle = v$.
\end{proof}

\section{Engine~II, One-Shot Dual Perturbation (Sharpening $11 \to 8$), and the $m$-Chain Lift}
\label{sec:08:engine-two}

Engine~I (\cref{prop:08:engine-one}) expresses the inert derivative valuation bound in terms of $F_{\mathrm{num}}(\ell, d, D_\ell)$, where the true degree profile $D_\ell$ differs slightly from the smooth uniform weight profile $W = (k - 1)w$ by a small signed perturbation
\begin{equation}
\label{eq:08:D-perturb}
  D_\ell \;=\; W + E - E', \qquad W = (k - 1) w, \qquad 0 \le E \le 4 w, \qquad 0 \le E' = 2\mathbf{1}_{\TT} w \le 2 w,
\end{equation}
arising from the top-set shifts $u \in \TT$ and the bihomogeneous jet weights of Engine~II (weight quantization and generic frame biorthogonalization on the Dwork power-series algebra).

To relate $F_{\mathrm{num}}(\ell, d, D_\ell)$ cleanly back to $F_{\mathrm{num}}(\ell, d, W)$, we first record the one-line incoming-flow complementary slackness estimate and then prove our \emph{eight-line one-shot dual test}, which eliminates the archive's two-step perturbation split and sharpens the explicit loss constant from $11 = 3 + (4 + 4)$ down to $8$!

\begin{lemma}[Incoming-flow complementary slackness]
\label{lem:08:slackness}
Let $\ell \ge 3$, $d \ge 2$ even, and let $\sigma_W \in \mathcal{D}_{\ell, d}$ be an optimal dual vector for $W = (k - 1)w$ ($w : \ZZ/d\ZZ \to \RR_{\ge 0}$, $k \ge 3$, $S := \sum_{s \in \ZZ/d\ZZ} w(s)$). Then
\begin{equation}
\label{eq:08:slackness}
  \sum_{\substack{s \in \ZZ/d\ZZ \\ \sigma_W(s) \ge 2/\ell}} w(s) \;\le\; \frac{1}{\ell}\, S.
\end{equation}
\end{lemma}

\begin{proof}
Let $x^* \in \RR_{\ge 0}^E$ be an optimal primal flow for $w$ ($A x^* \le w$, $\mathbf{1}^T x^* = \langle \sigma_W, w \rangle$ by \cref{thm:08:lp-duality}, after dividing $W = (k-1)w$ by $k-1 > 0$), and set $B := \{s \in \ZZ/d\ZZ : \sigma_W(s) \ge 2/\ell\}$. Expanding the duality gap \eqref{eq:08:weak-duality},
\[
  0 \;=\; \langle \sigma_W, w \rangle - \mathbf{1}^T x^*
  \;=\; \sum_{t \in \ZZ/d\ZZ} \sigma_W(t)\bigl(w(t) - (A x^*)(t)\bigr)
  \;+\; \sum_{(s, i) \in E} \bigl(\ell^i \sigma_W(s) + \sigma_W(s + i) - 1\bigr)\, x^*(s, i).
\]
Because every term in both sums is nonnegative, every term vanishes identically. For every $s \in B$:
\begin{itemize}
\item $\sigma_W(s) \ge 2/\ell > 0$ forces $(A x^*)(s) = w(s)$; and
\item for every $i \in \mathrm{oddLevels}(d)$, $\ell^i \sigma_W(s) + \sigma_W(s + i) - 1 \ge \ell \cdot \frac{2}{\ell} - 1 = 1 > 0$ forces $x^*(s, i) = 0$.
\end{itemize}
Consequently, for every $s \in B$ we have $w(s) = (A x^*)(s) = \sum_{i \in \mathrm{oddLevels}(d)} x^*(s - i, i)$ (purely incoming flow!). Summing over $s \in B$ and noting that each $x^*(s - i, i)$ contributes $\ell^i x^*(s - i, i) \ge \ell x^*(s - i, i)$ to $(A x^*)(s - i) \le w(s - i)$ gives
\[
  \sum_{s \in B} w(s) \;=\; \sum_{s \in B} \sum_{i} x^*(s - i, i)
  \;\le\; \frac{1}{\ell}\sum_{t \in \ZZ/d\ZZ}\sum_{i} \ell^i x^*(t, i)
  \;\le\; \frac{1}{\ell}\sum_{t \in \ZZ/d\ZZ} (A x^*)(t) \;\le\; \frac{1}{\ell}\, S. \qedhere
\]
\end{proof}

We now prove our eight-line one-shot dual test sharpening the archive's explicit odd-level gain constant from $11 = 3 + (4 + 4)$ down to $8$.

\newsimplification{Eight-Line One-Shot Dual Test Sharpening the Odd-Level Inert Gain Constant from $11$ to $8$}

\begin{proposition}[Sharpened odd-level perturbation bound (constant $8 < 11$)]
\label{prop:08:const-eight}
Let $\ell \ge 3$, $d \ge 2$ even, $k \ge 3$, $w : \ZZ/d\ZZ \to \RR_{\ge 0}$ with $S := \sum_{s \in \ZZ/d\ZZ} w(s)$, and $W := (k - 1)w$. Suppose $D_\ell = W + E - E' \ge 0$ with $0 \le E \le 4 w$ and $0 \le E' \le 2 w$ pointwise on $\ZZ/d\ZZ$. Then
\begin{equation}
\label{eq:08:const-eight}
  F_{\mathrm{num}}(\ell, d, D_\ell) \;\ge\; F_{\mathrm{num}}(\ell, d, W) - \frac{8}{\ell}\, S
  \;\ge\; F_{\mathrm{num}}(\ell, d, W) - \frac{11}{\ell}\, S.
\end{equation}
\end{proposition}
\leanref{src/Formalism/GaussJetBounds/OneLetterGaussJets/OneLetterGaussProofs/CompanionTheory.lean:38}{Formalism.DworkPowerSeriesWeights.OddLevelGainBounds.proposition\_six}

\begin{proof}
Let $\sigma_W \in \mathcal{D}_{\ell, d}$ be an optimal dual vector for $W$ (\cref{thm:08:lp-duality}), so $F_{\mathrm{num}}(\ell, d, W) = \sum_s W(s)(\gamma_\ell - \sigma_W(s))$. Testing $D_\ell$ directly against the \emph{same} feasible dual vector $\sigma_W$ via weak duality \eqref{eq:08:weak-duality} gives
\[
  F_{\mathrm{num}}(\ell, d, D_\ell)
  \;\ge\; \sum_{s \in \ZZ/d\ZZ} D_\ell(s)\bigl(\gamma_\ell - \sigma_W(s)\bigr)
  \;=\; F_{\mathrm{num}}(\ell, d, W) + \sum_{s \in \ZZ/d\ZZ} \bigl(E(s) - E'(s)\bigr)\bigl(\gamma_\ell - \sigma_W(s)\bigr).
\]
Recall from \eqref{eq:08:gamma-p} that $\frac{1}{\ell} < \gamma_\ell = \frac{\ell}{\ell^2 - 1} \le \frac{3}{2\ell} < \frac{2}{\ell}$ for all $\ell \ge 3$. Partition $\ZZ/d\ZZ$ into three disjoint regimes according to $\sigma_W(s) \in [0, 1]$:
\begin{enumerate}[label=\textup{(\arabic*)}]
\item $0 \le \sigma_W(s) \le \gamma_\ell$: then $\gamma_\ell - \sigma_W(s) \in [0, \gamma_\ell]$ and $E(s) - E'(s) \ge -2w(s)$, so $(E - E')(\gamma_\ell - \sigma_W) \ge -2\gamma_\ell w(s) \ge -\frac{3}{\ell} w(s) \ge -\frac{4}{\ell} w(s)$.
\item $\gamma_\ell < \sigma_W(s) < 2/\ell$: then $\sigma_W(s) - \gamma_\ell \in (0, \frac{1}{\ell})$ and $E(s) - E'(s) \le 4w(s)$, so $(E - E')(\gamma_\ell - \sigma_W) \ge -\frac{4}{\ell} w(s)$.
\item $\sigma_W(s) \ge 2/\ell$: then $\sigma_W(s) - \gamma_\ell \in (0, 1]$ and $E(s) - E'(s) \le 4w(s)$, so $(E - E')(\gamma_\ell - \sigma_W) \ge -4 w(s)$.
\end{enumerate}
Summing over all $s \in \ZZ/d\ZZ$ and invoking incoming-flow complementary slackness (\cref{lem:08:slackness}) yields
\[
  F_{\mathrm{num}}(\ell, d, D_\ell)
  \;\ge\; F_{\mathrm{num}}(\ell, d, W) - \frac{4}{\ell}\sum_{\sigma_W(s) < 2/\ell} w(s) - 4\sum_{\sigma_W(s) \ge 2/\ell} w(s)
  \;\ge\; F_{\mathrm{num}}(\ell, d, W) - \frac{4}{\ell} S - 4 \cdot \frac{1}{\ell} S
  \;=\; F_{\mathrm{num}}(\ell, d, W) - \frac{8}{\ell} S. \qedhere
\]
\end{proof}

We now state the full \emph{Inert-Prime Derivative Divisibility Theorem} combining Engine~I, Engine~II (biextension weight quantization and generic frame biorthogonalization), odd-orbit floor lifting, and the $m$-step chain lift across all tower layers $0 \le j \le m$.

\begin{theorem}[Inert-prime derivative divisibility and $m$-chain lift]
\label{thm:08:inert-deriv}
Let $\ell \ge 5$ be a level-free inert prime (\cref{def:08:level-free}), and consider the definite unitary period tower of ranks $N_j = k + j(t+1)$ ($0 \le j \le m$) evaluated at CM torsion points along the host CM type $\Phi$.
\begin{enumerate}[label=\textup{(\roman*)}]
\item \emph{(Engine~II: Weight quantization and generic frame biorthogonalization.)} In the Dwork power-series algebra attached to the height-$2$ biextension $G_{c,\sigma(c)}$ (\cref{thm:07:biextension-rigidity}), the formal-group jet weights satisfy the exact binomial quantization identity along one-letter Gauss jets; consequently, on any Frobenius orbit $\mathcal{O}$ of length $d$ with generic degree frame, there exists a biorthogonal dual basis over $\OO_{\CC_\ell}$ whose diagonal entries are controlled by the orbital LP value $F_{\mathrm{num}}(\ell, d, D_\ell)$ up to the universal perturbation bound \eqref{eq:08:const-eight} (constant $8 < 11$).
\item \emph{(Odd-orbit floor lift.)} When the Frobenius orbit $\mathcal{O}$ on $\Phi$ has odd length $d_{\mathrm{odd}}$ (so the full $\Frob_\ell$-orbit on $\Phi \cup c\Phi$ has even length $2d_{\mathrm{odd}}$), unfolding $w$ periodically to $\ZZ/(2d_{\mathrm{odd}})\ZZ$ preserves the dual polytope $\mathcal{D}_{\ell, 2d_{\mathrm{odd}}}$ and yields the exact same $-\frac{8}{\ell} S$ floor bound.
\item \emph{($m$-chain derivative divisibility package.)} Summing over all $\Frob_\ell$-orbits partitioning $\Phi$ and chaining across all $m$ steps $0 \le j \le m$ of the period tower yields, for every top index $u \in \TT$ and translate $u_0 \in A_0$, a lower bound on the local $\ell$-adic determinant valuation of the column-chain evaluation matrix $\mathcal{E}_\ell(u, u_0)$ matching the inert-prime contribution of $\ASK(D, N, a, u, w_{u_0})$ at $\ell$ up to an explicit orbital rounding error of size $O\bigl(k\, S_{\mathrm{orb}} / \ell + m d_{\mathrm{orb}}\bigr)$, ready for coarse-multiplicity extraction and global gluing in \cref{ch:09}.
\end{enumerate}
\end{theorem}
\leanref{src/Formalism/GaussJetBounds/GaussJetProofs/OneLetterRadiusIdentity.lean:48}{Formalism.DworkPowerSeriesWeights.weightQuantization}
\leanref{src/Formalism/GaussJetBounds/GaussJetProofs/FrameA.lean:287}{Formalism.DworkPowerSeriesWeights.theorem\_A\_frame}
\leanref{src/Formalism/SharpRestriction/Divisibility/OrbitPlaces.lean:362}{Formalism.MainNotions.OrbitPlaces.EntriesAtInert}

\begin{proof}
(i) For each one-letter Gauss jet attached to $\exp(G_{c,\sigma(c)}(X_s, X_{s+i}))$, the radius-of-convergence identity fixes the $\ell$-adic weight quantization of every monomial $X_s^a X_{s+i}^b$ identically in terms of the binomial sums of $\ell$-power digits. Biorthogonalizing the resulting polynomial jet frame over the DVR $\OO_{\CC_\ell}$ via Smith normal form (producing an invertible integral transition matrix modulo $\varpi$-powers determined by the primal flow LP \eqref{eq:08:flow-constraint}) expresses the local column-chain divisibility exponent directly in terms of $F_{\mathrm{num}}(\ell, d, D_\ell)$, where $D_\ell = (k-1)w + E - E'$ satisfies \eqref{eq:08:D-perturb}. Applying \cref{prop:08:const-eight} replaces $F_{\mathrm{num}}(\ell, d, D_\ell)$ by $F_{\mathrm{num}}(\ell, d, (k-1)w) - \frac{8}{\ell} S$.

(ii) When $d = 2 d_{\mathrm{odd}}$ is the full orbit length on $\Phi \cup c\Phi$, every $\sigma \in \mathcal{D}_{\ell, 2d_{\mathrm{odd}}}$ minimizes $\langle \sigma, W \rangle$ on the periodic extension of $w$ across $\ZZ/(2d_{\mathrm{odd}})\ZZ$; testing $D_\ell$ against $\sigma_W$ in \cref{prop:08:const-eight} requires only that $2d_{\mathrm{odd}}$ is even, yielding the identical $-\frac{8}{\ell} S$ loss.

(iii) Evaluating the exact LP main term $F_{\mathrm{num}}(\ell, d, (k-1)w)$ on the two-level thick/thin weight profile $w_{u_0}$ of \eqref{eq:03:two-level-weight}: at each orbit position $s$ where $a / \ell^i > w_{u_0}(\pi_u(i\Frob_\ell))$, the optimal dual vector extracts precisely $(a/\ell^i - w_{u_0}(\pi_u(i\Frob_\ell)))_+$ weighted by $2\log\ell$, reproducing the $\ell$-summand of $\ASK(D, N, a, u, w_{u_0})$ in \eqref{eq:03:ask-def}! Summing across all orbits partitioning $\Phi$ and chaining across the $m$ steps $j = 0, 1, \dots, m - 1$ accumulates the divisibility exponent additively across layers while keeping the boundary and perturbation loss bounded by $O(k S / \ell + m d)$, completing the proof.
\end{proof}

